\section{Introduction}
\label{sec:intro}

Modern software development demands constant movement across tools: a task
management board to plan work, a chat platform to coordinate with teammates,
and a code-hosting service to review pull requests and track issues. Each
transition carries a context-switching cost---the cognitive overhead of
re-establishing mental state in a new tool---that fragments attention and slows
progress~\cite{mark2008cost,czerwinski2004diary}. This fragmentation is
compounded when information lives in separate silos: a chat thread that
references a board card requires the developer to manually open two applications
and mentally merge their contents.

Integrated developer workspaces attempt to reduce this overhead by co-locating
surfaces. Yet the empirical record on whether integration alone suffices---and
what further benefit is conferred by embedding an AI agent that can observe and
act across those surfaces---remains thin. Prior work on AI-assisted programming
has focused heavily on code generation~\cite{copilot2022,chen2021codex}, leaving
the broader coordination and project-management layer largely unstudied.

A second open question concerns the role of \emph{context breadth}: how much
workspace state should an agent be allowed to observe? Giving an agent access
to more surfaces may improve its relevance, but it raises concerns about
information overuse, trust, and auditability. Designing studies to isolate this
variable has proven difficult because prior systems enforce context boundaries
in prompts---a mechanism that is opaque, brittle, and unverifiable after the
fact.

\textbf{This paper makes three contributions:}
\begin{enumerate}
  \item \textbf{\Weld system.} We describe \Weld, an open-source all-in-one
    developer workspace (Kanban board, team chat, GitHub integration) augmented
    with a study mode that provisions isolated per-participant workspaces and
    logs fine-grained behavioral telemetry~(\S\ref{sec:system}).

  \item \textbf{Code-enforced context scoping (\ContextGate).} We introduce the
    \ContextGate, a server-side module that assembles agent context from a
    condition-specific allowlist. Because scoping is enforced in code---not
    prompts---the manipulation is mechanically verifiable from audit
    logs~(\S\ref{sec:system}).

  \item \textbf{Empirical study.} We report a between-subjects user study
    ($N$~=~\emph{TBD}) crossing three agent-context conditions against three
    task types, with measures spanning task performance, cognitive load
    (\NASATLX), usability (\SUS), and cross-surface telemetry~(\S\ref{sec:results}).
\end{enumerate}

The remainder of the paper is structured as follows.
Section~\ref{sec:related} surveys related work.
Section~\ref{sec:system} describes \Weld and the study instrumentation.
Section~\ref{sec:study} presents the study design.
Section~\ref{sec:results} reports results.
Section~\ref{sec:discussion} discusses implications, and
Section~\ref{sec:conclusion} concludes.
